% Part: normal-modal-logic
% Chapter: axioms-systems
% Section: introduction

\documentclass[../../../include/open-logic-section]{subfiles}

\begin{document}

\olfileid[id]{nml}{axs}{int}

\olsection{Pengantar}

Kita mempunyai semantik bagi bahasa modal dasar dalam pengertian model modal,
serta pengertian bahwa !!a{formula} valid---benar pada semua dunia dalam semua
model---atau valid terhadap suatu kelas model atau kerangka---benar pada semua
dunia dalam semua model yang termasuk dalam kelas itu atau berbasis pada
kerangka tersebut. Logika biasanya menghubungkan pencirian semantik validitas
semacam ini dengan pengertian teori-bukti tentang !!{derivability}. Tujuannya
ialah mendefinisikan pengertian !!{derivability} dalam suatu sistem sedemikian
sehingga !!a{formula} !!{derivable} jika dan hanya jika formula itu valid.

Sistem !!{derivation} yang paling sederhana dan secara historis paling tua
ialah sistem !!{derivation} jenis Hilbert atau aksiomatik. Sistem
!!{derivation} jenis Hilbert bagi banyak logika modal relatif mudah dibangun:
sistem-sistem ini sederhana sebagai objek kajian metateoretis (misalnya, untuk
membuktikan kesahihan dan kelengkapan). Akan tetapi, sistem-sistem tersebut
jauh lebih sulit digunakan untuk membuktikan !!{formula}s daripada, misalnya,
sistem deduksi alami.

Dalam sistem !!{derivation} jenis Hilbert, !!a{derivation} dari !!a{formula}
merupakan barisan !!{formula}s yang berawal dari aksioma-aksioma tertentu,
melalui beberapa aturan inferensi, dan berakhir pada !!{formula} yang
bersangkutan. Karena kita menginginkan agar sistem !!{derivation} cocok dengan
semantik, kita harus menjamin bahwa himpunan formula yang !!{derivable} benar
dalam semua model (atau benar dalam semua model tempat semua aksioma benar).
Mula-mula kita akan mengisolasi beberapa sifat logika modal yang diperlukan
untuk tujuan tersebut: logika modal ``normal''. Bagi logika modal normal,
hanya ada dua aturan inferensi yang perlu diasumsikan: modus ponens dan
nesesitasi. Sebagai aksioma, kita mengambil semua instansiasi substitusi
tautologi dan, bergantung pada logika modal yang kita hadapi, sejumlah aksioma
modal. Bahkan jika kita hanya tertarik pada kelas semua model, kita juga harus
menghitung semua instansiasi substitusi dari~$\Ax{K}$
\iftag{notprvDiamond}{}{ dan~$\Ax{Dual}$} sebagai aksioma. Hal ini saja sudah
menghasilkan logika modal normal minimal~$\Log K$.

\begin{defn}
Aturan \emph{modus ponens} merupakan skema inferensi
\begin{prooftree}
\AxiomC{$!A$}
\AxiomC{$!A \lif !B$}
\RightLabel{\MP}
\BinaryInfC{$!B$}
\end{prooftree}
Kita mengatakan bahwa !!a{formula}~$!B$ \emph{mengikuti dari}
!!{formula}s~$!A$, $!C$ melalui modus ponens jika dan hanya jika
$!C\ident !A\lif !B$.
\end{defn}

\begin{defn}
Aturan \emph{nesesitasi} merupakan skema inferensi
\begin{prooftree}
\AxiomC{$!A$}
\RightLabel{\Nec}
\UnaryInfC{$\Box !A$}
\end{prooftree}
Kita mengatakan bahwa !!{formula}~$!B$ mengikuti dari !!{formula}~$!A$
melalui nesesitasi jika dan hanya jika $!B\ident\Box !A$.
\end{defn}

\begin{defn}
Suatu \emph{!!{derivation}} dari himpunan aksioma~$\Sigma$ merupakan barisan
!!{formula}s $!B_1$, $!B_2$, \dots, $!B_n$, dengan setiap $!B_i$ merupakan
salah satu dari yang berikut:
\begin{enumerate}
\item suatu instansiasi substitusi dari tautologi;
\item suatu instansiasi substitusi dari !!a{formula} dalam~$\Sigma$;
\item mengikuti dari dua !!{formula}s $!B_j$, $!B_k$ dengan $j$, $k<i$ melalui
  modus ponens; atau
\item mengikuti dari !!a{formula}~$!B_j$ dengan $j<i$ melalui nesesitasi.
\end{enumerate}
Jika terdapat !!a{derivation} demikian dengan $!B_n\ident !A$, kita mengatakan
bahwa $!A$ \emph{!!{derivable} dari~$\Sigma$}, dengan lambang
$\Sigma\Proves !A$.
\end{defn}

Dengan definisi ini, akan ternyata bahwa himpunan formula yang !!{derivable}
membentuk suatu logika modal normal, dan bahwa setiap !!{formula} yang
!!{derivable} benar dalam setiap model tempat setiap aksioma benar. Sifat
!!{derivation}s ini disebut \emph{kesahihan}. Konversnya, yakni
\emph{kelengkapan}, lebih sulit dibuktikan.

\end{document}
